\documentclass[11pt]{article}
\usepackage[margin=1in]{geometry}
\usepackage{amsmath,amssymb,amsthm}
\usepackage{booktabs}
\usepackage[numbers,sort&compress]{natbib}
\usepackage[expansion=false]{microtype}
\usepackage{xcolor}
\usepackage[colorlinks=true,linkcolor=blue!50!black,citecolor=blue!50!black,urlcolor=blue!50!black]{hyperref}
\hypersetup{pdftitle={A complete kernel-only Lean proof of the three-copy Werner endpoint for every dimension},pdfauthor={Tyler Duic}}
% Generated by scripts/generate_numbers.py from shipped evidence.
% Do not edit by hand; regenerate instead.
\newcommand{\nmodules}{8,199}
\newcommand{\nprep}{9}
\newcommand{\nlight}{7,829}
\newcommand{\nheavy}{145}
\newcommand{\nfinal}{216}
\newcommand{\nretries}{62}
\newcommand{\runid}{41356979845}
\newcommand{\walltime}{12 h 3 min}
\newcommand{\cpuhours}{18.0}
\newcommand{\peakmem}{6.40}
\newcommand{\peakmod}{F31.P576x1}
\newcommand{\nmodulesbefore}{8,110}
\newcommand{\runstart}{2026-11-01 22:26:47 UTC}
\newcommand{\runend}{2026-11-02 10:29:39 UTC}
\newcommand{\maxreldiff}{3e-15}
\newcommand{\nblocks}{1,000}
\newcommand{\nrows}{273}
\newcommand{\certbits}{2.11M}
\newcommand{\ncpu}{8}
\newcommand{\containergb}{9}
\newcommand{\certshrunk}{11.1}
\newcommand{\vonebits}{23.48M}
\newcommand{\certbytes}{1,638,079}
\newcommand{\dendelbits}{131}
\newcommand{\vonedendelbits}{5,941}


\newtheorem{theorem}{Theorem}
\newtheorem{lemma}[theorem]{Lemma}
\theoremstyle{remark}
\newtheorem{remark}[theorem]{Remark}

\newcommand{\C}{\mathbb{C}}
\newcommand{\Tr}{\operatorname{Tr}}

\title{A complete kernel-only Lean proof\\ of the three-copy Werner endpoint\\ for every dimension}
\author{Tyler Duic\\[2pt]\small Independent researcher, New York, NY\\\small ORCID \href{https://orcid.org/0009-0005-4675-3045}{0009-0005-4675-3045}\\\small\url{https://github.com/tduic}}
\date{November 2026}

\begin{document}

\maketitle

\begin{abstract}
For $C$ an operator on $(\C^d)^{\otimes 3}$ of rank at most $2$, let
$q(C) = \sum_{S \subseteq \{1,2,3\}} (-1/2)^{|S|} \lVert \Tr_S C \rVert_F^2$.
The inequality $q(C) \ge 0$ is the three-copy Werner endpoint.
I prove nothing new about Werner states in this note: the endpoint was
already proved for every $d \ge 3$ by Tomar~\cite{tomar2026}, whose Lean
formalization covers $3 \le d \le 8$ only.
What I contribute is a complete kernel-only Lean~4 proof of the endpoint
for every $d$, by an independent route: one exact rational certificate of
the $t = 1/d \to 0$ limit of a strengthened moment relaxation, carried to
every dimension by a lemma chain, with the certificate's exact identity
proved in the Lean kernel.
The final theorem depends only on the standard axioms
\texttt{propext}, \texttt{Classical.choice}, and \texttt{Quot.sound},
with no \texttt{sorry} and no \texttt{native\_decide}, and I had its
statement independently checked to confirm it says what I claim here.
\end{abstract}

\section{Introduction}

The Werner state $\rho_\alpha \propto I - \alpha F$ on $\C^d \otimes \C^d$,
where $F$ is the swap, is a central test case for distillability.
For $k$ copies, undistillability reduces to a quadratic-form inequality on
operators $C$ of rank at most $2$ on $(\C^d)^{\otimes 3}$ when $k = 3$.
Define, for $\alpha \in \mathbb{R}$,
\begin{equation}
q_\alpha(C) = \sum_{S \subseteq \{1,2,3\}} (-\alpha)^{|S|} \lVert \Tr_S C \rVert_F^2,
\qquad q = q_{1/2}.
\end{equation}
Then $\rho_\alpha$ is three-copy undistillable exactly when $q_\alpha \ge 0$
on rank-$\le 2$ operators $C$, and the continuation identity used in the
earlier work on this problem shows that
$q_\alpha \ge (1-2\alpha)(1-\alpha)^2 \lVert C \rVert^2$ on $[0, 1/2]$
holds in dimension $d$ exactly when the endpoint $q \ge 0$ holds there.
The endpoint is therefore the sharp three-copy statement.

Tomar~\cite{tomar2026} proves the endpoint for every $d \ge 3$,
by a base certificate polynomial in $1/d$ for $d \ge 9$ plus exact
certificates for smaller $d$.
That paper is a preprint, not peer reviewed, and its Lean formalization
covers $3 \le d \le 8$ only: the paper states that
``the case $d \ge 9$ is not formalized''.
I close that formalization gap, by a different certificate route, and I
claim nothing else as new.

\begin{theorem}\label{thm:main}
For every $d$ and every operator $C$ on $(\C^d)^{\otimes 3}$ of rank at
most $2$, $q(C) \ge 0$.
\end{theorem}

Theorem~\ref{thm:main} is Tomar's theorem, restated.
My contribution is Theorem~\ref{thm:lean} below:
the same statement, proved in Lean~4 by the kernel, for every $d$ at once.

\section{The moment relaxation and its $t = 0$ limit}

I prove Theorem~\ref{thm:main} in Lean through a strengthened moment
relaxation.
A rank-$\le 2$ operator $C$ lifts to moments built from its factors;
twirling over the relevant symmetries lands in a commutative algebra
spanned by diagram operators, and the question becomes whether a certain
moment functional is nonnegative on the physical moment set.
Replacing the orbit variables by trace coordinates $m$ makes
the relaxation $d$-free except for the site modules: the objective and
the equality rows are trace functionals, so they have $d$-free
coefficients in $m$, while only the positive-semidefinite blocks depend
on $d$, through site matrices $N_r(p)(t)$ with $t = 1/d$.

At $t = 0$ no site module degenerates: every module's limit matrices
span the full matrix algebra. An exact rational certificate
$X_b \succeq 0$ and multipliers $\lambda$ for the $t = 0$ relaxation
therefore suffices, by the lemma chain of Section~\ref{sec:chain},
to prove $q(C) \ge 0$ in every dimension.
The certificate I use is the compact v2 certificate shipped in this
folder. It passes the standalone exact verifier \texttt{verify\_t0.py}:
\nblocks{} blocks are positive definite.
It is $\certshrunk{}\times$ smaller in bits than the earlier v1 certificate
(\certbits{} bits versus \vonebits{} bits), and its correction
denominator is \dendelbits{} bits (v1: \vonedendelbits{} bits).
These are properties of the shipped certificate data, checked by exact
rational arithmetic outside Lean; Section~\ref{sec:lean} is what moves
the same identity into the Lean kernel.

\section{The lemma chain}\label{sec:chain}

The chain is an abstract theorem in Lean, \texttt{endpoint\_of\_t0\_cert}:
given a $t = 0$ certificate, the endpoint follows in every dimension.
Its steps are:

\begin{enumerate}
\item \emph{Embedding.} A physical point of dimension $e$ embeds into
$d \ge \max(e, 4)$ without changing $m$ or $q$, because the permutation
operators commute with the embedding and partial traces commute with it.
\item \emph{Validity at dimension $d$.} The twirled coordinates
satisfy every homogeneous row, the trace row equals $1$, the objective
equals $q(C)$, and every word block and coupled block is positive
semidefinite, for any real tensor-space module basis.
\item \emph{$d$-freeness in $m$.} Every row and the objective are trace
functionals, so their $m$-forms are exactly their $d$-pipeline forms
composed with the Weingarten change of variables. The qubit factor
causes no gap, since the conjugated moments lie in the range of the
qubit Gram matrix.
\item \emph{Blocks at $t = 1/d$.} For $d \ge 4$, the blocks built from
$N_r(p)(1/d)$ are positive multiples of compressions of the true blocks:
the site data are positive multiples of Gram-form matrices, the
one-site Gram matrix at the identity is positive definite, and the
denominator $(1-d^{-2})(1-4d^{-2})(1-9d^{-2})$ is positive.
Hence the blocks are positive semidefinite for every physical $m$ of
dimension $e \le d$.
\item \emph{Continuity.} The site matrices are polynomial in $t$, so the
blocks converge as $t = 1/d \to 0$, and the positive-semidefinite cone
is closed. Every physical $m$, of every dimension, therefore lies in
the $t = 0$ feasible set.
\item \emph{Pairing.} At a physical $m$,
$q(C) = \sum_b \langle X_b, L_b(0)(m) \rangle + \lambda_{\mathrm{trace}}$,
a sum of nonnegative pairings plus a nonnegative trace multiplier.
Hence $q(C) \ge 0$.
\end{enumerate}

I proved each step in Lean in the \texttt{AllD} layers, with no
\texttt{sorry} and only the standard axioms.
The chain is generic: it does not depend on the particular certificate,
only on the existence of one satisfying the $t = 0$ certificate
conditions. The remaining work is therefore exactly one thing:
proving the certificate identity for the explicit v2 data, in the kernel.

\section{Layer D: the certificate identity in the kernel}\label{sec:lean}

The one hypothesis of \texttt{endpoint\_of\_t0\_cert} that is not generic
is the exact identity for the explicit certificate data, \texttt{IdentityT0}:
for every orbit-invariant $\mu$, the objective equals the sum over blocks
of $\langle X_b, L_b(\mu) \rangle$ plus the row terms.
I prove this identity in the Lean kernel, in three parts.

First, the \emph{bridge} rewrites each block pairing as a sum over family
parts, with a common integer scale, so that the identity becomes a
vanishing statement for family totals.
Second, the \emph{glue} reduces that statement to \texttt{FamZero} for
the three families, i.e.\ to literal-free family totals, using only the
orbit invariances of $\mu$; coupled block parts fold through a
kernel-checked full-rank re-presentation.
Third, the \emph{compute} checks the family totals in the kernel, in
\nmodules{} generated identity modules pinned by SHA-256 in
\texttt{identity\_manifest.json}, together with the hashes of the four
certificate files (\texttt{blocks.jsonl}, \texttt{multipliers.json},
\texttt{qubit.json}, \texttt{sites.json}).
A rep-budget split was needed to fit the largest claim modules under
the process memory limit; it raised the module count from \nmodulesbefore{} to
\nmodules{} without changing any family total.

The complete build ran as run \runid{}, in one
\ncpu{}-CPU, \containergb{}~GB container, from a clean checkout, starting
\runstart{} and finishing \runend{}
(\walltime{} wall time, \cpuhours{} CPU-hours summed over the modules).
All \nmodules{} modules kernel-checked with return code $0$
(prep \nprep{}, light \nlight{}, heavy \nheavy{}, final \nfinal{});
\nretries{} attempts of the parallel light pass were killed for memory and passed
in the serial retry; the run summary records no Lean error anywhere.
Peak memory was \peakmem{}~GB for one module.

\begin{theorem}\label{thm:lean}
In Lean~4, with Mathlib, the theorem
\texttt{LayerDT0.werner3\_endpoint\_all\_d} states that for every
$d : \mathbb{N}$ and every
$C : \mathrm{Matrix}\,(\mathrm{Fin}\,3 \to \mathrm{Fin}\,d)\,
(\mathrm{Fin}\,3 \to \mathrm{Fin}\,d)\,\C$ with $C.\mathrm{rank} \le 2$,
$0 \le \mathrm{qfun}\,C$, where \texttt{qfun} is the Lean definition of $q$.
It depends only on the axioms \texttt{propext},
\texttt{Classical.choice}, and \texttt{Quot.sound}.
\end{theorem}

The same axiom list holds for \texttt{LayerDT0.identityT0}.
I scanned every shipped Lean file for \texttt{sorry}, \texttt{admit},
\texttt{axiom}, \texttt{native\_decide}, \texttt{implemented\_by},
\texttt{extern}, \texttt{unsafe}, and \texttt{opaque} outside comments;
the scan found none (see \texttt{tests/test\_lean\_scan.py} in this folder).
An independent agent that had not seen the work producing the proof
checked the final theorem's statement and rated it faithful:
\texttt{qfun} is proved equal in Lean to an independently written $q$
(the six proper partial traces written out on index triples, Mathlib's
\texttt{Matrix.trace} for the full trace, \texttt{Complex.normSq} for
squared entries), rank is Mathlib's \texttt{Matrix.rank} with no
coercions, no shadowing declarations or extra axioms were found in the
import chain, $d = 1$ gives $q = \lVert c \rVert^2/8$ in Lean as expected,
and an independent NumPy comparison on random complex $C$ gave maximum
relative difference approximately $\maxreldiff{}$.
The proof term itself was out of scope for that check; the kernel build
above is what checks it.

\section{Prior art, and what I do not claim}

Tomar~\cite{tomar2026} proves Theorem~\ref{thm:main} for every $d \ge 3$.
What is new here is Theorem~\ref{thm:lean}: as far as I have found, the
first complete kernel-only Lean proof of the endpoint for every $d$,
closing the formalization gap at $d \ge 9$ left by Tomar~\cite{tomar2026}.
If an earlier complete kernel-only Lean proof for every $d$ exists,
I ask to be told, and I will correct this claim.
Nothing here concerns any other number of copies, and my claim is the
all-$d$ proof only.

\section*{Acknowledgements}

I did this work with substantial help from AI models
(Anthropic's Claude): searching, writing code and Lean proofs, and
running the independent reviews described above.
I directed the work and I am responsible for every claim;
that is why everything here is backed by code, data, tests, and
machine-checked proofs that anyone can rerun from this folder alone.

\bibliographystyle{plainnat}
\bibliography{references}

\end{document}
